\documentclass{article}
\usepackage[T1]{fontenc}
\usepackage{lmodern}
\usepackage{graphicx}
\usepackage{amsmath,amssymb}
\usepackage[a4paper,margin=2cm]{geometry}
\usepackage[absolute,overlay]{textpos}
\setlength{\TPHorizModule}{1mm}
\setlength{\TPVertModule}{1mm}
\usepackage[indonesian]{babel}
\usepackage{hyperref}
\setlength{\emergencystretch}{2em}
% unit: Halaman 82

\begin{document}

% === Halaman 82 (python/native) ===

Kriptanalisis Playfair Cipher

• Karena ada 26 huruf abjad, maka terdapat 26 x 26 = 677 bigram, sehingga 

identifikasi bigram individual lebih sukar.

• Sayangnya ukuran poligram di dalam Playfair cipher tidak cukup besar, 

hanya dua huruf sehingga Playfair cipher  tetap tidak aman. 

• Meskipun Playfair cipher sulit dipecahkan dengan analisis frekuensi relatif 

huruf-huruf, namun ia dapat dipecahkan dengan analisis frekuensi 

pasangan huruf. 

• Dalam Bahasa Inggris kita bisa mempunyai  frekuensi kemunculan 

pasangan huruf, misalnya pasangan huruf TH dan HE paling sering 

muncul. 

82

\end{document}
